\documentclass[11pt]{article}
\usepackage{parskip}
\usepackage{color}

\title{Pirogies}
\author{Nathaniel Zuk}

\begin{document}
\maketitle

I got this recipe from the Pirogies book (simple dough; potatoes ruski filling), but I think it needs some adjustments.  The techniques are particularly important.

\textbf{Ingredients:} (this makes about 60 pirogies)

Dough:

\begin{itemize}
\item Flour (~5 cups)
\item Warm water (2 cups)
\item Eggs (2)
\item Salt (2 tsp)
\end{itemize}

Combine 4 cups flour with the salt.  Mix in the eggs, then add the water 1/2 cup at a time.  As the texture becomes more liquidy when fully combined, add more flour 1/2 cup at a time.  When the dough begins to hold together, turn it out on a floured surface and knead until the dough is smooth and not too sticky (it will still be a bit sticky when it's done), about 10 min.  Wrap in plastic wrap and put in the fridge for at least 10 min.

Filling:

\begin{itemize}
\item Potatoes, mashed (5 lbs)
\item Onions, finely chopped (2) \textcolor{red}{(273 g)}
\item Ricotta (8 oz) \textcolor{red}{(5 oz)}
\item Eggs (2)
\item Salt (2 tsp)
\end{itemize}

Sautee the onions until golden brown.  Mix the eggs and ricotta into the potatoes.  Add the onions and salt and mix.

\textit{To make the pirogies:}

Work with one half of the dough at a time.  Roll out the dough to 1/8 inch thick.  Cut 3-inch diameter circles out of the dough.  Combine the dough and repeat rolling it out and cutting it until you finish with the dough (this should produce ~30 circles for each half).

While holding the circle, place a spoonful of filling in the middle.  Fold the dough in half around the filling, pressing the edges closed.  Let rest on a greased or floured cookie sheet \textcolor{red}{(it sticks to parchment paper!)}.  

As you continue, the circles that are resting while you work will dry out on the side face up.  When you get to dough circles that are dried out, flip them over and put the filling on the stickier side of the dough.  This makes it easier to seal.

When they are done briefly blanch the pirogies before storing. \textcolor{red}{(This wasn't part of the original directions, but when I tried the directions they provided, flouring the pirogies, they really stuck together.  I found that the boiled pirogies do not stick together, so blanching might be the best thing to do before freezing them.)}

\end{document}